\documentclass[11pt]{article}
\usepackage[a4paper,margin=22mm]{geometry}
\usepackage{fontspec}
\setmainfont{LinLibertine_R.otf}[BoldFont=LinLibertine_RB.otf,ItalicFont=LinLibertine_RI.otf,BoldItalicFont=LinLibertine_RBI.otf]
\usepackage{xcolor}
\usepackage{hyperref}
\hypersetup{colorlinks=true,urlcolor=blue}
\setlength{\parindent}{0pt}
\setlength{\parskip}{6pt}
\sloppy
\emergencystretch=2em
\begin{document}
\section*{Historical Britannica: extracted source}
Encyclopaedia Britannica - Second edition, Volume 4, D-F - EB.4\par
{\small Source row 510. Parser: nvidia/nemotron-3.5-parse.\par}
{\small This is model-extracted text, not a corrected transcription. OCR errors, reading order and region alignment require human review.\par}
\hrule\medskip
DRAW BRIDGE
Plate XCVI
DRAWING.
\par

{\scriptsize\color{gray} Region region-0; provider order 0.\par}
\medskip

7
8
A
\par

{\scriptsize\color{gray} Region region-1; provider order 1.\par}
\medskip
\newpage\section*{Source record}
bri-0510-56402624e6eb\par
Image SHA-256: \texttt{56402624e6eb7001ba19a34a2311f864}\par\texttt{1ec61116297d5491c5576485659d18ce}\par
Dataset revision: \texttt{309f765b1ecc83e0a7f7}\texttt{762e44b4a927a7555443}\par
National Library of Scotland, Encyclopaedia Britannica Illustrated (CC0).\par
\url{https://huggingface.co/datasets/NationalLibraryOfScotland/encyclopaedia_britannica_illustrated/viewer/default/train?row=510}
\end{document}
